\section{材料与方法}
\label{sec:methods}
\subsection{研究设计与证据来源}
研究方案和正式协议在正式结果产生前保存于本地。其形式属于预先规定的本地方案，不是公开注册平台上的预注册。工程预试验使用一个种子、三个重复刺激和一个恢复分支；通过后没有依据曲线调整正式输入频率、脉冲宽度、STD 参数、主读出或判据。

正式实验比较两种模型、两个刺激起始间隔和五个输入实现。每条序列含 12 次重复刺激，随后从同一训练末态派生两个独立恢复分支。所有条件都使用同一个连接组。研究流程见图\ref{fig:design}；完整协议为 \path{habituation/protocols/main_v1.json}。

\begin{figure}[!htb]
\centering
\small\textbf{A. 模型与读出：解剖结构保留，候选慢机制局部改变传递}\par\medskip
\resizebox{0.97\linewidth}{!}{%
\begin{tikzpicture}[>=Latex,font=\footnotesize,
 box/.style={draw=Ink!65,rounded corners=2pt,align=center,minimum height=1.0cm,minimum width=2.4cm,fill=Ink!3}]
\node[box] (input) at (0,0) {冻结随机输入\\逐次相同回放};
\node[box] (jon) at (3.1,0) {70 个 JO-CE\\机械感觉细胞};
\node[box,minimum width=3.1cm] (brain) at (6.7,0) {保留完整连接网络\\127,400 个神经元};
\node[box,minimum width=3.1cm] (readout) at (10.8,0) {aBN1（主要读出）\\aDN1/aDN2（次要读出）};
\draw[->,thick] (input)--(jon);
\draw[->,thick] (jon)--(brain);
\draw[->,thick] (brain)--(readout);
\node[align=center,text=Blue] (std) at (4.8,-1.5)
{M0：固定传递强度\\M1：JO-CE 兴奋性输出边\\引入可恢复资源状态};
\draw[->,Blue] (std.north)--(4.5,-0.1);
\node[box,dashed,fill=white,minimum height=0.85cm] (body) at (10.8,-1.65) {完整身体与梳理动作\\本轮未建模};
\draw[->,dashed] (readout)--(body);
\end{tikzpicture}}

\medskip\textbf{B. 连续训练与独立恢复分支（示意，不按时间比例绘制）}\par\medskip
\resizebox{0.95\linewidth}{!}{%
\begin{tikzpicture}[>=Latex,font=\footnotesize,
 pulse/.style={draw=Blue,fill=Blue!9,rounded corners=2pt,align=center,minimum width=1.1cm,minimum height=0.8cm},
 state/.style={draw=Ink,align=center,minimum width=1.8cm,minimum height=0.9cm}]
\node[pulse] (p1) at (0,0) {刺激 1};
\node[pulse] (p2) at (1.7,0) {刺激 2};
\node (ellipsis) at (3.1,0) {$\cdots$};
\node[pulse] (p12) at (4.5,0) {刺激 12};
\node[state,fill=Ink!5] (snap) at (7,0) {同一训练末态\\保存快照};
\draw[->] (p1)--node[above] {$T$}(p2);
\draw[->] (p2)--(ellipsis);\draw[->] (ellipsis)--(p12);\draw[->] (p12)--(snap);
\node[state,minimum width=3.0cm] (r2) at (11,0.8) {休息 2 s\\$\rightarrow$ 探测 A};
\node[state,minimum width=3.0cm] (r10) at (11,-0.8) {休息 10 s\\$\rightarrow$ 探测 A};
\draw[->] (snap.east)--++(0.5,0)|-(r2.west);
\draw[->] (snap.east)--++(0.5,0)|-(r10.west);
\node[align=center] at (2.6,-1.05) {刺激宽度 200 ms；响应窗口 300 ms\\$T=0.5$ s 或 2 s；序列内部不重置};
\end{tikzpicture}}
\caption{实验逻辑示意。A 中箭头表示输入和测量流程，不是完整解剖回路图；JO-CE 与被监测细胞均属于同一全脑网络，图中方框不是互斥细胞集合。M1 改变 JO-CE 的全部选中兴奋性输出边，不限于直接连到 aBN1 的边。B 中两个探测分支独立恢复同一训练末态，因此短休息探测不会污染长休息结果。}
\label{fig:design}
\end{figure}


\subsection{连接图谱与模型边界}
固定上游实现为 Shiu 等人的公开代码\citep{shiu2024}，提交为
\begin{center}\small\texttt{91bdd1e7dcf193f3e7ca5a8933497fcef63b7960}\end{center}
采用 FlyWire v630 处理后的细胞和连接表：127,400 个神经元、14,687,178 条聚合有向边，对应 52,793,639 个解剖突触。没有为降低成本而裁剪子图；仿真期间保留所有细胞和静态连接。这里的“全脑”沿用该模型范围，不包括虚拟身体、完整腹神经索或全部生理过程。

70 个 JO-CE 输入 ID 从固定版本 Notebook 的字面量清单提取，不执行整个 Notebook。主要读出为 aBN1，次要描述性读出为 aDN1/aDN2；精确标识列于附录表\ref{tab:ids}。本文并未对原论文的全部实验进行完整复现。

\subsection{M0：原始固定连接 LIF 模型}
对非不应期中的神经元 $j$，状态方程为
\begin{align}
\tau_m\frac{\mathrm{d}v_j}{\mathrm{d}t}&=v_0-v_j+g_j,\label{eq:lif}\\
\tau_g\frac{\mathrm{d}g_j}{\mathrm{d}t}&=-g_j.\label{eq:g}
\end{align}
$v_j$ 为膜电位，$g_j$ 为上游实现中的有效突触驱动，\textbf{单位同样是电压，不是电导}。当 $v_j>v_{\mathrm{th}}$ 时发出脉冲，将 $v_j$ 置为 $v_{\mathrm{rst}}$、$g_j$ 置零，并进入不应期。上游实现把上述两个微分方程都标记为在不应期内停止更新；本研究没有更改该约定。数值积分使用 Brian2 的线性状态更新方法及 Cython 后端\citep{stimberg2019}。

来自细胞 $i$ 的脉冲经过固定传递延迟后，令
\begin{equation}
g_j\leftarrow g_j+w_{ij},\qquad
w_{ij}=n_{ij}s_iw_{\mathrm{syn}},\label{eq:weight}
\end{equation}
其中 $n_{ij}$ 为聚合突触数量，$s_i$ 是上游传递符号约定，$w_{\mathrm{syn}}$ 为全局尺度。正式实验保持上游静态参数不变，见表\ref{tab:params}。固定连接的含义是没有显式、随刺激历史变化的突触效能规则；它不表示 $v$ 和 $g$ 没有历史。

\begin{table}[tbp]
\centering\small
\caption{固定模型参数与本研究新增参数。M1 的两个 STD 参数为预先选定的演示设置，并非 JO-CE 生理拟合结果。}
\label{tab:params}
\begin{tabularx}{\linewidth}{lllY}
\toprule
类别 & 符号 & 数值 & 含义\\
\midrule
静态 LIF & $v_0,v_{\mathrm{rst}}$ & $-52\mV$ & 静息与重置电位\\
& $v_{\mathrm{th}}$ & $-45\mV$ & 发放阈值\\
& $\tau_m$ & $20\ms$ & 膜时间常数\\
& $\tau_g$ & $5\ms$ & 突触驱动衰减时间\\
& $t_{\mathrm{rfc}}$ & $2.2\ms$ & 普通细胞不应期；被驱动输入细胞设为零\\
& $t_{\mathrm{dly}}$ & $1.8\ms$ & 突触传递延迟\\
& $w_{\mathrm{syn}}$ & $0.275\mV$ & 每个解剖突触的全局权重尺度\\
数值/输入 & $\Delta t$ & $0.1\ms$ & 积分及输入事件网格\\
& $f$ & $150\,\mathrm{Hz}$ & 刺激期间的等效单细胞输入率\\
新增 STD & $U$ & $0.01$ & 每次到达事件的资源消耗比例\\
& $\tau_{\mathrm{rec}}$ & $2\secunit$ & 资源恢复时间常数\\
\bottomrule
\end{tabularx}
\end{table}

\begin{primer}{怎样理解时间常数}
时间常数不是“多久后绝对归零”，而是指数变化的速度尺度。一个单纯衰减的量经过一个时间常数后约剩原来的 37\%。原模型的显式单细胞时间常数为毫秒级，而新增资源状态的恢复尺度为秒级。不过，不能仅凭这个差别排除网络产生慢模式的可能，所以仍需要实际运行 M0。
\end{primer}

\subsection{M1：JO-CE 兴奋性输出的局部短时突触抑制}
\label{sec:std}
M1 只替换以 JO-CE 为起点、静态权重为正的 4,764 条聚合输出边。对每条选中连接引入无量纲可用资源 $x_{ij}\in[0,1]$，初始为 1。事件之间满足
\begin{equation}
\frac{\mathrm{d}x_{ij}}{\mathrm{d}t}=\frac{1-x_{ij}}{\tau_{\mathrm{rec}}}.
\label{eq:resource}
\end{equation}
当一个突触前脉冲到达时，先用到达前资源传递，再消耗资源：
\begin{align}
g_j&\leftarrow g_j+w_{ij}x_{ij}^{-},\\
x_{ij}^{+}&=(1-U)x_{ij}^{-}.\label{eq:deplete}
\end{align}
这里的“抑制”指短时传递效能下降（depression），\textbf{不是把这些兴奋性连接变成 GABA 能抑制性连接}。该机制也不同于“加强一组抑制性神经元”的竞争假说。

恢复采用事件驱动的精确指数更新。对于相邻到达事件，先把上一次事件后的状态推进到当前时刻，再执行式\eqref{eq:deplete}。仿真实现中将对应原始静态边的权重置零，另建同起点、同终点、同延迟的动态边，避免双重传递；其他连接不变。它没有新增解剖通路，也没有按刺激序号人为乘一个越来越小的包络。

\begin{primer}{暂时用掉的“资源”是什么}
可以把 $x$ 想成一个使用后需要补充的缓冲器：$x=1$ 时传递完整，$x=0.5$ 时同一个脉冲在这条边上只产生一半有效驱动。真实生物中可能有多个过程导致类似现象；本研究的 $x$ 是抽象状态，不是直接测到的囊泡数，也没有验证其一定对应 JO-CE 的真实生理过程。
\end{primer}

“一个恢复时间尺度”不等于“全脑只有一个状态变量”。当前实现按边存储资源，所有边共享 $U$ 和 $\tau_{\mathrm{rec}}$。在相同初态、相同延迟与参数下，同一突触前细胞对应的这些边具有相同资源更新历史，理论上可以共享状态；本轮没有实施该优化。两项共享参数没有通过数据拟合，4,764 条边也不是 4,764 个被独立训练的参数。

\subsection{输入：冻结随机实现与严格配对}
每个输入神经元在刺激持续的 200 ms 内，在每个 $0.1\ms$ 时间格独立抽取一次 Bernoulli 事件，概率为
\begin{equation}
p=f\Delta t=150\times 10^{-4}=0.015.
\end{equation}
这是小步长的离散稀疏随机驱动，不是严格连续时间 Poisson 过程。单次刺激的理论输入事件均值为 $70\times150\times0.2=2,100$。随机种子为 1701--1705。

每个种子的事件阵列只生成一次，然后在该种子的所有重复刺激、两个模型、两个间隔和恢复探测中回放。因此，同一序列内不会因“下一次碰巧更少的输入”而产生假递减。跨种子改变输入实现，以检查结果对有限输入变化的稳健性。

事件通过强数值驱动注入输入神经元的膜电位，幅度为上游的 $250w_{\mathrm{syn}}=68.75\mV$；与上游一样，将这些被驱动细胞的不应期设为零。这个注入是促使模型细胞发放的数值接口，\textbf{不是对真实机械刺激的生理幅度校准}。注入事件仍需经过 LIF 阈值与重置调度，不能未经检查就假定每个事件必然一对一变成感觉脉冲。分析器因此比较实际感觉神经元发放的逐事件序列，而不只比较注入计数。

\subsection{连续状态刺激与独立恢复分支}
\label{sec:protocol}
第 $k$ 次刺激（$k=1,\ldots,12$）的起始时刻为 $(k-1)T$，$T\in\{0.5,2\}\secunit$。刺激持续 0.2 s；响应在起始后的 0.3 s 窗口内计数。因此，两个协议的\emph{刺激结束到下次刺激开始}静默间隔分别为 0.3 s 和 1.8 s，不能把 $T$ 本身误写为静默间隔。

\textbf{整个重复序列内不重建网络，不清空膜电位、突触驱动、不应期或待传递事件。}仅在切换独立模型条件、输入实现或刺激间隔时恢复初态。完整训练段结束时刻是 $11T+0.3\secunit$，短/长间隔分别为 5.8 s 和 22.3 s。

训练末态保存后，分支 A 等待 2 s 再给一次原刺激，分支 B 重新恢复\emph{同一个末态}、等待 10 s 后再给原刺激。休息从末次 300 ms 响应窗口结束开始计，而不是从末次刺激结束计；所以相对于末次刺激结束，两个探测间隔实际为 2.1 s 和 10.1 s。这样的分支设计避免了 2 s 探测本身改变 10 s 结果。

\subsection{读出、效应量与预先规定的筛查}
\label{sec:metrics}
对每个种子、模型和刺激间隔，令 $R_k$ 为第 $k$ 次刺激后 300 ms 内 aBN1 的脉冲数。定义
\begin{equation}
I=R_1,\qquad E=\frac{R_1+R_2+R_3}{3},\qquad
L=\frac{R_{10}+R_{11}+R_{12}}{3},\qquad
D=1-\frac{L}{E}.\label{eq:D}
\end{equation}
其中 $I$ 是首次响应，$E$ 为早期均值，$L$ 为末期均值，$D$ 为递减比例。$E=0$ 时不定义该比值，不能把无初始响应当作完全习惯化。本文先逐种子计算 $D$，再平均；它通常不等于先汇总 $L,E$ 后再作比值。

H1 的实用性筛查阈值为 $D\geq0.30$。在 H1 通过的前提下，H2 筛查要求 10 s 探测 $P_{10}$ 同时满足
\begin{equation}
P_{10}-L\geq0.20E,\qquad P_{10}\geq0.80I.\label{eq:H2}
\end{equation}
平台筛查要求第 7--9 次与第 10--12 次响应均值之差不超过 $0.10E$。这些是本项目预先规定的描述性门槛，不是文献对习惯化的普遍数值定义，也不是统计显著性检验。

对恢复另报告 $P_t/I$。看过部分结果后增加的探索性诊断是 $(P_t-L)/(I-L)$，仅在 $I>L$ 时定义；它衡量“已经损失的响应恢复了多少”，与“回到首次水平的多少”不同。该探索性指标没有被事后加入 H4 的通过判据。

\subsection{统计单位、工程对照与软件证据}
五个种子是同一连接组上的五种输入实现，\textbf{不是五只动物}。一条序列内的 12 个响应也不是 12 个独立生物样本。结果以逐种子值、均值、极值范围和筛查通过数描述；没有构造生物总体置信区间或宣称显著性。

正式实验前，M1-zero 将 $U$ 设为零，检查替换突触的实现是否与 M0 等价；三个 pilot 模型分别通过训练末态恢复重放检查。正式实验用独立分析入口回读原始事件，校验文件哈希、细胞 ID 精度、时间网格、重复事件、资源边界、刺激时序和窗口计数，再验证 20 条序列的完整覆盖。

使用 Python 3.10.21、Brian2 2.5.1、NumPy 1.22.3、Pandas 1.4.3、Cython 0.29.33 与 PyArrow 17.0.0；主仿真单进程运行，限制原生数学库线程。完整环境沿用既有独立复现环境，未修改系统 Python 或上游源码。软件、参数和数据的固定方式见附录\ref{app:repro}。
\FloatBarrier
